% !TeX root = ../../Book3.tex
\section{单连通 Riemann 曲面}
\label{b3:sec:1201}

\chapref{b3:ch:11}已经构造每张连通 Riemann 曲面的全纯万有覆盖，并证明底曲面是源曲面对甲板群的商。现在还缺少一个解析问题：单连通的源曲面能够采用什么全局复坐标？本节先研究球面、平面、圆盘三个具体模型，算清前两者允许的自由甲板群，并比较圆盘的两种循环商。\secref{b3:sec:1202}的统一化定理才负责证明这些模型穷尽全部可能。

\subsection{球面}
\label{b3:subsec:120101}

Riemann 球面 \(\widehat{\mathbb C\kern0pt}\) 使用有限坐标 \(z\) 与无穷远坐标 \(w=1/z\)，重叠处的转换全纯且导数非零。两个坐标闭圆盘 \(|z|\le1\)、\(|w|\le1\) 覆盖全曲面，所以它紧。

\ProseParagraphBreak
球面还单连通，但不能以“每条连续闭路都漏掉一个点”作为理由，连续路径的像可能非常复杂。对一条给定闭路，用有限小坐标圆盘覆盖其像，再细分参数，使每段落在一个坐标圆盘内。在坐标中把该段直线替换，凸性给出端点固定同伦。替换后的路径只由有限条圆弧或线段组成，可以选一个未被经过的球面点 \(q\)。删去 \(q\) 后的球面经一个 Möbius 坐标等价于平面，替换路径便可在其中收缩。原闭路也随之可收缩。这里先作有限局部替换，才使用去掉一点的平面坐标。

\ProseParagraphBreak
球面上的亚纯函数已经由\cref{b3:thm:03-sphere-rational}证明为有理函数。自同构比一般有理函数更严格：它在每个点都有全纯局部逆，因而不能出现重复覆盖。下面用唯一极点及其阶数把这一限制落实。

\Needspace{6\baselineskip}
\begin{proposition}[球面不能有非平凡自由甲板群]\label{b3:prop:12-sphere-free}
每个双全纯自同构 \(T:\widehat{\mathbb C\kern0pt}\to\widehat{\mathbb C\kern0pt}\) 都是 Möbius 变换
\[
 T(z)=\frac{az+b}{cz+d},\quad a,b,c,d\in\mathbb C,
 \quad ad-bc\ne0.
\]
若 \(T\) 非恒等，则它至少有一个球面不动点。因此，若 \(X\) 为连通 Riemann 曲面，\(p:\widehat{\mathbb C\kern0pt}\to X\) 为全纯万有覆盖，则 \(p\) 是双全纯映射。
\end{proposition}
\begin{proof}
令 \(q=T^{-1}(\infty)\)。由于 \(T\) 单射，\(q\) 是它唯一的极点；在源坐标及目标无穷远坐标中，\(T\) 有全纯局部逆，所以这个极点是一阶的。选 Möbius 变换 \(M\)，使 \(M(\infty)=q\)：若 \(q\) 有限可取 \(M(z)=q+1/z\)，否则取恒等映射。于是 \(R=T\circ M\) 只在无穷远有一个一阶极点，在有限平面全纯。

将球面全纯映射看作亚纯函数，\cref{b3:thm:03-sphere-rational}给 \(R=P/Q\)。约去公因子后，分母不能有有限零点，故由代数基本定理为常数；因此 \(R\) 是多项式。它在无穷远的极点阶数就是多项式次数，所以次数为1。\(T=R\circ M^{-1}\) 因而是 Möbius 变换，且行列式非零。

若 \(c=0\)，\(T\) 固定无穷远。若 \(c\ne0\)，有限不动点满足
\[
 cz^2+(d-a)z-b=0.
\]
代数基本定理保证至少一个根。这个根不可能同时使 \(cz+d=0\)，否则原方程还会给 \(az+b=0\)，违背 \(ad-bc\ne0\)。所以它确实是有限不动点。

最后，甲板变换全纯，而\cref{b3:prop:11-deck-free}说明非恒等甲板变换不能固定源点，故 \(p\) 的甲板群只有恒等元。由于 \(p\) 是万有覆盖，\cref{b3:prop:11-universal-deck}保证甲板群在每条纤维上传递，所以每个纤维只有一点。覆盖又局部双全纯，因而这个全局双射及其逆都全纯。
\end{proof}
这里起作用的是“所有球面自同构都有不动点”与“甲板作用自由”的结合。紧性本身不能排除非平凡覆盖：后面的平移格商也紧，却由整个平面覆盖。另一方面，球面上的全局复值全纯函数必恒常，因为模在紧球面上取得最大值，局部最大模原理及连通性迫使函数恒常；有理函数的丰富性来自允许取值 \(\infty\)。

\subsection{平面}
\label{b3:subsec:120102}

平面 \(\mathbb C\) 凸，所以单连通；它没有球面的紧性。\Cref{b3:thm:05-entire-univalent}已经证明，单射整函数恰为 \(z\mapsto az+b\)、\(a\ne0\)。因此平面自同构的自由作用受到一个直接限制：若 \(a\ne1\)，点 \(b/(1-a)\) 被固定；非恒等甲板变换只能是非零平移。

\ProseParagraphBreak
多个平移能够组成什么群，取决于它们能否无限接近恒等变换。覆盖要求每个小层与所有非恒等甲板像不交，因而排除了任意短的非零平移。以下证明把这项局部限制转成平面上的整数格。

\Needspace{6\baselineskip}
\begin{proposition}[平面的离散平移群]\label{b3:prop:12-translation-lattice}
设 \(X\) 为连通 Riemann 曲面，\(p:\mathbb C\to X\) 为全纯万有覆盖。则所有甲板变换均为平移，且平移向量群
\[
 \Lambda=\{\,\lambda\in\mathbb C\mid
              (z\mapsto z+\lambda)\in\operatorname{Deck}(p)\,\}
\]
恰为下列三种形式之一：\(\{0\}\)、\(\mathbb Zu\)、\(\mathbb Zu+\mathbb Zv\)。第二种中 \(u\ne0\)；第三种中 \(u,v\) 在 \(\mathbb R\) 上线性独立。反过来，每个所列加法群的平移作用都自由且适当不连续，其商具有以 \(\mathbb C\) 为源的全纯万有覆盖。
\end{proposition}
\begin{proof}
由平面自同构分类及甲板作用自由性，每个甲板变换均为平移，向量在加法与取负下封闭。由\cref{b3:prop:11-deck-proper}，0有邻域 \(W\)，与其所有非恒等甲板像不交。选 \(\varepsilon>0\)，使 \(D(0,\varepsilon)\subset W\)。若非零 \(\lambda\in\Lambda\) 满足 \(|\lambda|<\varepsilon\)，则 \(\lambda\) 同时属于 \(W\) 与 \(W+\lambda\)，矛盾。因此每个非零向量的模都至少为 \(\varepsilon\)。

更重要的是，任意两个不同向量之差仍是非零群元，所以
\[
 |\lambda-\mu|\ge\varepsilon
 \quad(\lambda,\mu\in\Lambda,\ \lambda\ne\mu).
\]
用有限个半径小于 \(\varepsilon/2\) 的圆盘覆盖任一闭有界圆盘，每个小圆盘至多含一个群元，故 \(\Lambda\) 与任何有界集的交都是有限集。这一步使用了群的差封闭性，而不仅是逐点孤立。

若 \(\Lambda\ne\{0\}\)，先任选非零 \(\lambda_0\)。集合
\(\{\lambda\in\Lambda:0<|\lambda|\le|\lambda_0|\}\)
有限非空，从中取模最小的 \(u\)，便是全群最短非零向量。若 \(tu\in\Lambda\)、\(t\in\mathbb R\)，选整数 \(n\)，使 \(|t-n|\le1/2\)。向量 \((t-n)u\) 仍在群中，且模不超过 \(|u|/2\)，最短性迫使它为零。因此
\[
 \Lambda\cap\mathbb Ru=\mathbb Zu.
\]
若全部向量都在这条实直线上，得到第二种形式。

否则以 \(u\) 归一化坐标，写
\[
 \frac wu=a(w)+\mathrm i b(w),\quad a(w),b(w)\in\mathbb R.
\]
取 \(w_0\in\Lambda\)，经必要的变号使 \(b_0=b(w_0)>0\)。减去整数倍 \(u\) 可以把任何向量的横坐标归约到 \([0,1)\)，而不改变高度 \(b(w)\)。于是集合
\[
 E=\{\,w\in\Lambda\mid0\le a(w)<1,\quad0<b(w)\le b_0\,\}
\]
非空，且包含在一个有界平行四边形中，所以有限。取 \(v\in E\)，使 \(h=b(v)\) 最小。若全群还有正高度小于 \(h\) 的向量，先归约其横坐标，它也会落在 \(E\)，与最小性矛盾。故 \(h\) 是整个群的最小正高度；这同时证明正高度有正下界并且下界取得。

对任意 \(w\in\Lambda\)，取 \(n=\lfloor b(w)/h\rfloor\)，则
\[
 0\le b(w-nv)<h.
\]
最小正高度使这个高度只能为零。因此 \(w-nv\in\Lambda\cap\mathbb Ru=\mathbb Zu\)，得到 \(\Lambda=\mathbb Zu+\mathbb Zv\)。因为 \(b(v)>0\)，两个向量在实数域上线性独立。

反过来，所列群中的整数坐标使它们与有界集只有有限交点；秩二情形使用实线性同构 \((s,t)\mapsto su+tv\)，把有界集拉回有界集。非零平移没有固定点，故作用自由。若紧集 \(K,L\subset\mathbb C\) 满足 \((K+\lambda)\cap L\ne\varnothing\)，则 \(\lambda\in L-K\)，而 \(L-K\) 紧，只有有限个这样的群元。由\cref{b3:prop:11-proper-action-quotient}，商成为 Riemann 曲面，商投影是全纯覆盖；源平面单连通，故为万有覆盖。
\end{proof}
三种形式中的秩，分别指0、1、2个独立的整数生成元。两个实方向已被最短向量与最小高度用尽，不能在保持离散性的同时再加入第三个整数独立方向。商曲面是什么，则还须把坐标写出来。

\ProseParagraphBreak
秩0时没有不同点被识别，商就是平面。秩1时令
\[
 q_u(z)=\exp\left(\frac{2\pi\mathrm i z}{u}\right).
\]
指数的周期说明 \(q_u(z)=q_u(w)\) 当且仅当 \(z-w\in\mathbb Zu\)。每个非零复数都有对数，所以 \(q_u\) 满射到 \(\mathbb C^*\)；其导数处处非零。因此它在 \(\mathbb C/\mathbb Zu\) 上诱导一个双射，且在每个小商坐标中双全纯，得到
\[
 \mathbb C/\mathbb Zu\simeq\mathbb C^*.
\]
这项结论包含全部纤维和局部逆的核验，而不只是写出一个周期函数。

\ProseParagraphBreak
秩2时，每个 \(z\in\mathbb C\) 唯一写成 \(z=su+tv\)、\(s,t\in\mathbb R\)。分别减去 \(s,t\) 的整数部分，表明闭平行四边形
\[
 P=\{\,su+tv\mid0\le s\le1,\ 0\le t\le1\,\}
\]
映满商曲面，故商紧。上下两边及左右两边分别按平移识别。更明确地，记 \(\mathbb S^1=\{\zeta:|\zeta|=1\}\)，映射
\[
 [su+tv]\longmapsto
 \bigl(e^{2\pi\mathrm i s},e^{2\pi\mathrm i t}\bigr)
\]
是商到 \(\mathbb S^1\times\mathbb S^1\) 的连续双射；源紧而目标 Hausdorff，所以它是同胚。这给出了环面的拓扑，而复结构由平面小圆盘的商坐标给出。这样的 Riemann 曲面称为\term[fu-huanmian]{复环面}[complex torus]。实坐标表示说明其拓扑，并未把任意实线性变换宣布为全纯坐标变换。

\subsection{圆盘}
\label{b3:subsec:120103}

圆盘 \(\mathbb D\) 也凸、单连通且非紧。它与平面有相同的开圆盘拓扑，却有不同的全纯函数约束：恒等函数是非恒常有界全纯函数，而平面的有界整函数由 Liouville 定理只能常值。圆盘的自同构则通过 Cayley 变换成为上半平面的实 Möbius 变换。

\ProseParagraphBreak
由\cref{b3:prop:11-mobius-types}，一个没有内部固定点的非恒等圆盘自同构，在半平面坐标中属于抛物型或双曲型。对它生成的循环群，必要时更换生成元或缩放坐标，可以分别使用
\[
 T(w)=w+1,\quad S_a(w)=aw,\quad a>1
\]
作为模型。\subsecref{b3:subsec:110401}已经给出全部纤维及局部逆的计算，所得商为
\[
 \begin{aligned}
 \mathbb H/\langle T\rangle
 &\simeq\mathbb D\setminus\{0\},\\
 \mathbb H/\langle S_a\rangle
 &\simeq A\left(e^{-2\pi^2/\log a},1\right).
 \end{aligned}
\]
前者的商坐标是 \(e^{2\pi\mathrm i w}\)；后者在 \(0<\operatorname{Im}\log w<\pi\) 的分支中使用 \(\exp(2\pi\mathrm i\log w/\log a)\)。这些映射分别把整周期与伸缩轨道识别成同一点。

\ProseParagraphBreak
两商的基本群都为 \(\mathbb Z\)：它们的源单连通，甲板群又是无限循环群，可用\cref{b3:prop:11-fundamental-deck-action}识别。它们的万有覆盖也都等价于圆盘，却仍有不同的共形结构。双曲型循环商的模量为
\[
 \frac1{2\pi}\log\left(e^{2\pi^2/\log a}\right)
 =\frac{\pi}{\log a}.
\]
例如 \(a=e^\pi\) 给模量1，\(a=e^{2\pi}\) 给模量 \(1/2\)，由\cref{b3:prop:08-modulus-invariance}，这两个圆环不双全纯等价。去心盘的两端连接族具有无穷极值长度，\secref{b3:sec:0803}已用趋向穿孔的子圆环密度证明这一点；所以它也不等价于任何有限模量圆环。

\ProseParagraphBreak
平面的秩1商 \(\mathbb C^*\) 同样有基本群 \(\mathbb Z\)，但也不能与去心盘等价。否则将一个到去心盘的双全纯映射与指数 \(\mathbb C\to\mathbb C^*\) 复合，便得到非恒常有界整函数，违反 Liouville。由此，抽象基本群既不能决定万有覆盖的共形模型，也不能决定给定模型上甲板群的具体作用。

\subsection{三种基本模型}
\label{b3:subsec:120104}

以上讨论从具体模型出发，尚未把任意单连通曲面放进其中。关键区别在于全局坐标是否已经存在：Riemann 曲面的定义只给一组相容的局部坐标，Riemann 映照定理则以一个已经位于平面中的区域为出发点。下面把可以直接使用该定理的范围单独说清。

\Needspace{6\baselineskip}
\begin{proposition}[已嵌入球面的单连通曲面]\label{b3:prop:12-embedded-models}
设 \(X\) 为连通、单连通 Riemann 曲面。若存在全纯单射 \(F:X\to\widehat{\mathbb C\kern0pt}\)，则 \(X\) 双全纯等价于 \(\widehat{\mathbb C\kern0pt}\)、\(\mathbb C\)、\(\mathbb D\) 中恰好一种。
\end{proposition}
\begin{proof}
在源与目标的局部坐标中，单射全纯函数由\cref{b3:cor:01-local-injectivity}导数非零，所以 \(F\) 在每点局部双全纯。它的像 \(U\) 开，且局部逆与单射性拼成 \(F:X\to U\) 的全纯逆，故 \(U\) 单连通。

若 \(U\) 为整个球面，结论已得。否则取 \(a\notin U\)，用 Möbius 变换把 \(a\) 送到无穷远。像于是成为一个单连通平面区域 \(V\)。若 \(V=\mathbb C\)，得到平面模型；若 \(V\subsetneq\mathbb C\)，\cref{b3:thm:05-riemann}给 \(V\) 到圆盘的双全纯映射。

最后核对唯一性。球面紧，平面与圆盘不紧，故球面不能与后二者双全纯。若平面与圆盘双全纯，就存在非恒常有界整函数，与\cref{b3:thm:02-liouville}矛盾。因此三模型两两不等价。
\end{proof}
这项命题已完成“有全局单射坐标以后”的分类，却没有产生这个坐标。即使找到一个非恒常有界全纯函数，也还须证明它单射；圆盘上的 \(z\mapsto z^2\) 就表明有界性与非常值性不够。对紧单连通曲面，全局坐标还必须允许取无穷远值，不能只寻找复值全纯函数。

\ProseParagraphBreak
统一化的深层内容，正是对任意抽象单连通 Riemann 曲面构造这样的全局模型，而不是把已有平面区域再分类一次。有了覆盖模型，球面不动点、离散平移群与循环商的计算就能帮助辨认底曲面的形状。
